\section{Rational identities for a mixed cubic Hurwitz tail}
\label{mtt:section}
The preceding tail calculus suggests asking when a polynomially weighted
product has a rational ordinary sum after its divergent polynomial is
subtracted. The mixed-tail continuation~\cite{mtt:report} answers this
for the product of orders $(2,3)$ and proposes $(2,2,3)$ as a next case.
Here three explicit linear conditions give an infinite family for that
case. A finite difference certificate proves every member. The first
nonzero member has a degree-ten weight and value $-4037/108$.

Throughout this section put
\[
 A_n=\zeta(2,n+1),\qquad B_n=\zeta(3,n+1),\qquad x_n=n+\tfrac12.
\]
The only tail identities needed for the certificate are
\begin{equation}\label{mtt:eq:shifts}
 A_{n+1}=A_n-(n+1)^{-2},\qquad
 B_{n+1}=B_n-(n+1)^{-3}.
\end{equation}
For a polynomial $f$, let $\mathcal S f$ denote its unique polynomial
sum with $(\mathcal S f)(0)=0$ and
$(\mathcal S f)(n)-(\mathcal S f)(n-1)=f(n)$.
Thus, with Bernoulli polynomials in their usual convention,
\begin{equation}\label{mtt:eq:power-sum}
 (\mathcal S n^j)(n)
 =\frac{B_{j+1}(n+1)-B_{j+1}(1)}{j+1}\qquad(j\ge0).
\end{equation}
The variable in $\mathcal S n^j$ specifies the polynomial being summed.

\subsection{Three conditions and a finite certificate}
Let $P(n)=\sum_{d\ge5,\ d\ {\rm odd}}p_dn^d$ have rational coefficients,
and define
\begin{align}
 C&=\mathcal S(P/n^3),&
 D&=\mathcal S(2P/n^2),\label{mtt:eq:CD}\\
 F&=\mathcal S(-P/n^4+D/n^2).\label{mtt:eq:F}
\end{align}
These are polynomials. In particular $D$ is divisible by $n^2$, and the
odd powers in $-P/n^4+D/n^2$ cancel. Define three rational linear
functionals on $P$ by
\begin{equation}\label{mtt:eq:conditions}
 \lambda_0(P)=\sum_dp_dB_{d-3},\qquad
 \lambda_1(P)=\sum_ddp_dB_{d-3},\qquad
 \lambda_2(P)=F'(0)-\frac{p_5}{2}.
\end{equation}
The derivative in the last condition is evaluated from the finite
polynomial sum~\eqref{mtt:eq:F}; no infinite series is involved.

\begin{theorem}[Polynomial certificate space]\label{mtt:thm:certificate}
If $\lambda_0(P)=\lambda_1(P)=\lambda_2(P)=0$, then
\begin{align}
 E&=\mathcal S(2C/n^2-2P/n^5+D/n^3),\label{mtt:eq:E}\\
 G&=-P/n^7+C/n^4+D/n^5-E/n^2-F/n^3\label{mtt:eq:G}
\end{align}
are polynomials. Set
\begin{align}
 W_P(x)&=P(x+\tfrac12)-P(x-\tfrac12),\label{mtt:eq:weight}\\
 V_n&=P(n)A_n^2B_n+C(n)A_n^2+D(n)A_nB_n
             +E(n)A_n+F(n)B_n.\label{mtt:eq:V}
\end{align}
Then $V_0=0$ and
\begin{equation}\label{mtt:eq:finite}
 V_{n+1}-V_n=W_P(x_n)A_n^2B_n+G(n+1).
\end{equation}
Conversely, within the ansatz~\eqref{mtt:eq:V}, with all coefficient
polynomials zero at $0$ and $P\in n^5\mathbb Q[n^2]$, a polynomial
remainder in~\eqref{mtt:eq:finite} is possible exactly when these three
conditions hold.

On $\operatorname{span}_{\mathbb Q}\{n^5,n^7,\ldots,n^{2r+3}\}$ the
certificate space has dimension $\max(0,r-3)$. For $r=4$ it is generated by
\begin{equation}\label{mtt:eq:firstP}
 P(n)=35n^{11}+126n^9+75n^7+4n^5.
\end{equation}
These are statements about polynomial certificates. They do not assert
arithmetic independence of the periods occurring outside this space, or
classify every possible rational evaluation of a mixed-tail sum.
\end{theorem}
\begin{proof}
Write $C_j=[n^j]C$, and similarly for $D,E,F$. Faulhaber's formula gives
\[
 C_1=\lambda_0,\quad C_2=p_5/2,\quad
 D_1=0,\quad D_2=\lambda_1-2\lambda_0,\quad D_3=p_5.
\]
Consequently $2C/n^2-2P/n^5+D/n^3$ is an odd Laurent polynomial
whose only possible negative power is $\lambda_1/n$.
When $\lambda_1=0$, its polynomial sum defines $E$.
The input defining $F$ is even, as noted above. Since
$B_j(1)=0$ for odd $j>1$, the linear coefficient of the sum of an
odd polynomial is half its linear coefficient. The quadratic coefficient
of the sum of an even polynomial is half its quadratic coefficient.
Applied to the two inputs, this yields
\[
 E_1=C_3+D_4/2,\qquad F_2=D_4/2.
\]
The possible negative powers of $G$ are therefore precisely
\begin{equation}\label{mtt:eq:obstruction}
 [G]_{<0}
   =(\lambda_1-\lambda_0)n^{-3}-\lambda_2n^{-2};
\end{equation}
the $n^{-1}$ coefficient is $C_3+D_4-E_1-F_2=0$.
There are no lower powers. This proves polynomiality under the three
conditions.

Substituting~\eqref{mtt:eq:shifts} into~\eqref{mtt:eq:V} now proves
\eqref{mtt:eq:finite} by collecting the coefficients of
$A_n^2B_n,A_n^2,A_nB_n,A_n,B_n,1$. This also proves the converse:
the coefficients of $A_n^2,A_nB_n,B_n,A_n$ successively force
$C,D,F,E$, uniquely by their zero values at $0$.
The $A_n$ coefficient requires $\lambda_1=0$, and the constant
coefficient~\eqref{mtt:eq:obstruction} then requires
$\lambda_0=\lambda_2=0$.

The three functionals have the following matrix on $n^5,n^7,n^9$:
\[
 \begin{pmatrix}
  1/6&-1/30&1/42\\
  5/6&-7/30&3/14\\
  1/12&-7/180&23/630
 \end{pmatrix},
 \qquad \det=\frac1{12600}.
\]
Thus their rank is three as soon as $r\ge3$, and their first one or two
columns are independent. The dimension formula follows in every degree.
Substitution of~\eqref{mtt:eq:firstP} in the three conditions verifies
the stated generator.
\end{proof}

\subsection{Ordinary sums and rational endpoint values}
For an admissible $P$ from Theorem~\ref{mtt:thm:certificate}, let $Q_P(x)$
be the polynomial part of the asymptotic expansion of
$W_P(x)A_n^2B_n$ in $x=x_n$. It is explicit: multiply the finite portions
needed from
\begin{align}
 A_n&\sim\sum_{j\ge0}B_{2j}(\tfrac12)x_n^{-2j-1},\label{mtt:eq:centeredA}\\
 B_n&\sim\frac12\sum_{j\ge0}(2j+1)B_{2j}(\tfrac12)x_n^{-2j-2}.
 \label{mtt:eq:centeredB}
\end{align}
Only finitely many coefficients contribute to $Q_P$.
Both $W_P$ and $Q_P$ are even polynomials. The bracket below consequently
decays as $O(x_n^{-2})$.

\begin{corollary}[A rational ordinary-sum family]\label{mtt:cor:sum}
Every such $P$ gives the absolutely convergent identity
\begin{equation}\label{mtt:eq:ordinary}
 \sum_{n=0}^{\infty}
 \bigl(W_P(x_n)A_n^2B_n-Q_P(x_n)\bigr)
 =\operatorname{CT}_{N\to\infty}V_N\in\mathbb Q.
\end{equation}
Here $\operatorname{CT}$ is the coefficient of $N^0$ in the ordinary
large-$N$ expansion of~\eqref{mtt:eq:V}. It is a finite Bernoulli
calculation, used to evaluate an ordinary convergent sum.
\end{corollary}
\begin{proof}
Put
\[
 H_P=\mathcal S\bigl(G(n)+Q_P(n-\tfrac12)\bigr),\qquad U_n=V_n-H_P(n).
\]
Then $U_0=0$ and~\eqref{mtt:eq:finite} becomes
\[
 U_{n+1}-U_n=W_P(x_n)A_n^2B_n-Q_P(x_n).
\]
Finite-order Euler--Maclaurin expansions justify
\eqref{mtt:eq:centeredA}--\eqref{mtt:eq:centeredB} with remainders
to any required order. The centered parity proves absolute convergence,
so the telescoping partial sum tends to $\lim U_N$. Uncentered
Euler--Maclaurin expansions of $A_N$ and $B_N$ contain only rational
coefficients and integer powers of $N^{-1}$. Multiplication by the
fixed polynomials gives an expansion of $U_N$ of the same kind.
Since it has a finite limit, its positive powers vanish and its limit
is its constant coefficient. Finally $H_P(0)=0$, so the polynomial
$H_P(N)$ has no constant term. Thus
$\lim U_N=\operatorname{CT}V_N$, a rational number.
\end{proof}

\subsection{The first explicit identity}
For~\eqref{mtt:eq:firstP}, the weight and subtraction are
\begin{align}
 W(x)={}&\frac1{1024}\bigl(
 394240x^{10}+2639616x^8+4281984x^6\notag\\
 &\hspace{37mm}+1893344x^4+191316x^2+1995\bigr),
 \label{mtt:eq:explicitW}\\
 Q(x)={}&\frac{385}{2}x^6+\frac{3626}{3}x^4
                   +\frac{461503}{288}x^2+\frac{59471}{180}.
 \label{mtt:eq:explicitQ}
\end{align}
They give the ordinary identity
\begin{equation}\label{mtt:eq:explicit}
 \boxed{\displaystyle
 \sum_{n=0}^{\infty}
 \left(W(n+\tfrac12)\zeta(2,n+1)^2\zeta(3,n+1)
                    -Q(n+\tfrac12)\right)=-\frac{4037}{108}.}
\end{equation}
For a directly reproducible certificate the remaining polynomials are
\begin{align*}
 C(n)={}&\frac{n^2(n+1)^2(2n+1)}{18}
              (35n^4+70n^3+127n^2+92n+36),\\
 D(n)={}&n^3(n+1)^3(7n^4+14n^3+21n^2+14n+4),\\
 E(n)={}&\frac{n(n+1)}{216}
 (399n^6+1197n^5+6665n^4+11335n^3\\
 &\hspace{40mm}+16036n^2+10568n+5640),\\
 F(n)={}&\frac{n(n+1)(2n+1)}{18}
 (7n^6+21n^5+115n^4+195n^3\\
 &\hspace{40mm}+220n^2+126n+36),\\
 G(n)={}&-\frac{567n^6+7682n^4+27443n^2+16148}{216}.
\end{align*}
The correction in the proof of Corollary~\ref{mtt:cor:sum} is
\begin{align*}
 H_P(n)=\frac{n}{6480}\bigl(&175770n^6-8505n^5+1199985n^4-115230n^3\\
                         &+1989410n^2-411645n+860926\bigr).
\end{align*}
Thus the partial sum through $n=N-1$ is exactly $V_N-H_P(N)$.
The uncentered expansions
\[
 A_N\sim N^{-1}-\tfrac12N^{-2}
                 +\sum_{j\ge1}B_{2j}N^{-2j-1},\qquad
 B_N\sim\tfrac12N^{-2}-\tfrac12N^{-3}
                 +\tfrac12\sum_{j\ge1}(2j+1)B_{2j}N^{-2j-2}
\]
give $\operatorname{CT}V_N=-4037/108$, proving~\eqref{mtt:eq:explicit}.
All displayed asymptotic series are used only to finite order.

The experiment that found~\eqref{mtt:eq:firstP} was an exact rational
nullspace search for the three obstruction coefficients. The proof above
explains its rank in every degree. Independent verification constructs
nine polynomial certificates through degree $27$, checks the finite
difference identity symbolically, and recovers the subtraction and
endpoint by separate centered and uncentered expansions. Ordinary partial
sums at $70$ and $100$ digits agree with their finite certificates; four
deliberately corrupted identities are rejected. These finite checks
support the implementation, while the proofs establish the infinite
family. The analogous $(2,3,4)$ classification and the broader arithmetic
classification outside this certificate ansatz remain open.
